\documentclass[a4paper,11pt,twoside]{book}
\usepackage[T1]{fontenc}
\usepackage[utf8]{inputenc}
\usepackage{amsmath,amssymb,amsthm}
\usepackage{graphicx}
\usepackage{hyperref}
\usepackage{cleveref}

\newtheorem{theorem}{Theorem}[chapter]
\newtheorem{lemma}[theorem]{Lemma}
\newtheorem{definition}{Definition}[chapter]

\title{Foundations of Homotopy Type Theory}
\author{Dana Wallace}
\date{2026}

\begin{document}

\frontmatter
\maketitle
\tableofcontents
\listoffigures

\chapter{Preface}
This thesis examines the connections between homotopy theory and dependent type theory, with emphasis on univalence.

\mainmatter

\chapter{Preliminaries}
\label{chap:prelim}

\section{Type Theory Basics}
A judgement $a : A$ asserts that term $a$ has type $A$.

\begin{definition}
A type $A$ is called contractible when there exists an element $a : A$ such that for every $x : A$ we have $x = a$.
\end{definition}

\begin{theorem}[Univalence]
For types $A, B$ the canonical map $(A = B) \to (A \simeq B)$ is itself an equivalence.
\end{theorem}

\begin{proof}
See \Cref{chap:univalence} for the full argument.
\end{proof}

\chapter{Homotopical Structure}

\section{Paths and Loops}
Identity types $\text{Id}_A(x, y)$ correspond to paths in the space interpretation.

\section{Higher Inductive Types}
The circle $S^1$ is presented by a point and a loop.

\chapter{Univalence}
\label{chap:univalence}

The chapter develops the main technical machinery.

\backmatter

\chapter{Appendix A: Notation}
A brief reference for the notation used throughout.

\end{document}
